% figures/w03-ssa.svg — 각의 뺄셈은 각이 아니다. ssa 가 무엇을 고치는가.
%
%   bash _tools/tikz2svg.sh figures/src/w03-ssa.tex figures/w03-ssa.svg
%
% 이 파일은 손으로 쓰지 않는다 — _tools/w03_ssa.m 이 계산해서 낸다.
% 그 스크립트가 각·호의 방향·톱니 구간을 전부 assert 로 확인한 뒤에 쓴다.
%
% 회전 방향의 근거 (CLAUDE.md 4 규칙 6-2):
%   화면각 = 90 - 나침반각 이므로 나침반각이 커지면 화면각이 작아진다.
%   따라서 psi 가 커지는 것(우현)은 화면에서 시계방향이다.
%   raw = -340 -> 좌현이므로 화면 반시계, ssa = +20 -> 우현이므로 화면 시계.
\documentclass[border=5pt]{standalone}
% gnc-style.tex — 이 볼트의 TikZ 그림이 공유하는 스타일.
%
% 저널 그림 규약을 스타일 하나로 굳혀 둔다. 그림마다 선 굵기나 화살촉을
% 다시 정하지 않는다 — 그것이 그림이 제각각으로 보이는 원인이다.
%
% 쓰는 법:  % gnc-style.tex — 이 볼트의 TikZ 그림이 공유하는 스타일.
%
% 저널 그림 규약을 스타일 하나로 굳혀 둔다. 그림마다 선 굵기나 화살촉을
% 다시 정하지 않는다 — 그것이 그림이 제각각으로 보이는 원인이다.
%
% 쓰는 법:  % gnc-style.tex — 이 볼트의 TikZ 그림이 공유하는 스타일.
%
% 저널 그림 규약을 스타일 하나로 굳혀 둔다. 그림마다 선 굵기나 화살촉을
% 다시 정하지 않는다 — 그것이 그림이 제각각으로 보이는 원인이다.
%
% 쓰는 법:  \input{gnc-style}  (figures/src/ 안에서)

\usepackage{tikz}
\usepackage{amsmath}
\usepackage{xcolor}
\usetikzlibrary{arrows.meta, positioning, calc, fit, backgrounds, decorations.pathmorphing}

% ---- 색 --------------------------------------------------------------------
% 색은 강조 하나에만 쓴다. 나머지는 검정.
\definecolor{ink}{HTML}{1A1A1A}
\definecolor{accent}{HTML}{7C3AED}
\definecolor{warn}{HTML}{B91C1C}
\definecolor{soft}{HTML}{666666}

% ---- 선과 화살촉 -----------------------------------------------------------
\tikzset{
  % 신호선. 화살촉은 작고 뾰족한 것 하나뿐이다.
  sig/.style   = {ink, line width=0.5pt, -{Stealth[length=4pt, width=3pt]}},
  wire/.style  = {ink, line width=0.5pt},
  % 강조 경로 — 파선으로 구분한다. 흑백 인쇄에서도 살아야 하기 때문이다.
  asig/.style  = {accent, line width=0.5pt, densely dashed,
                  -{Stealth[length=4pt, width=3pt]}},
  awire/.style = {accent, line width=0.5pt, densely dashed},
  % 블록. 흰 바탕, 직각, 검은 테두리.
  blk/.style   = {draw=ink, line width=0.55pt, fill=white,
                  minimum width=17mm, minimum height=8mm, inner sep=2pt, font=\small},
  % 합산점
  sum/.style   = {draw=ink, line width=0.55pt, fill=white, circle,
                  inner sep=0pt, minimum size=4.6mm, font=\scriptsize},
  asum/.style  = {draw=accent, line width=0.55pt, fill=white, circle,
                  inner sep=0pt, minimum size=4.6mm, font=\scriptsize},
  % 분기점
  dot/.style   = {ink, fill=ink, circle, inner sep=0pt, minimum size=1.5mm},
  % 신호 이름 — 선 위에 작게
  sname/.style = {font=\small, inner sep=1.5pt},
  % 그림 안의 짧은 설명. 그림 안의 글은 최소로 둔다
  note/.style  = {font=\scriptsize, soft, inner sep=1.5pt},
  anote/.style = {font=\scriptsize, accent, inner sep=1.5pt},
  % 축
  axis/.style  = {ink, line width=0.5pt},
  tick/.style  = {font=\scriptsize, soft},
}

% 캡션 한 줄 — 그림 아래 가는 선과 함께
\newcommand{\gnccaption}[2]{%
  \node[anchor=north west, text width=#1, align=left, font=\scriptsize, ink]
        at (current bounding box.south west) {\vspace{2pt}#2};}
  (figures/src/ 안에서)

\usepackage{tikz}
\usepackage{amsmath}
\usepackage{xcolor}
\usetikzlibrary{arrows.meta, positioning, calc, fit, backgrounds, decorations.pathmorphing}

% ---- 색 --------------------------------------------------------------------
% 색은 강조 하나에만 쓴다. 나머지는 검정.
\definecolor{ink}{HTML}{1A1A1A}
\definecolor{accent}{HTML}{7C3AED}
\definecolor{warn}{HTML}{B91C1C}
\definecolor{soft}{HTML}{666666}

% ---- 선과 화살촉 -----------------------------------------------------------
\tikzset{
  % 신호선. 화살촉은 작고 뾰족한 것 하나뿐이다.
  sig/.style   = {ink, line width=0.5pt, -{Stealth[length=4pt, width=3pt]}},
  wire/.style  = {ink, line width=0.5pt},
  % 강조 경로 — 파선으로 구분한다. 흑백 인쇄에서도 살아야 하기 때문이다.
  asig/.style  = {accent, line width=0.5pt, densely dashed,
                  -{Stealth[length=4pt, width=3pt]}},
  awire/.style = {accent, line width=0.5pt, densely dashed},
  % 블록. 흰 바탕, 직각, 검은 테두리.
  blk/.style   = {draw=ink, line width=0.55pt, fill=white,
                  minimum width=17mm, minimum height=8mm, inner sep=2pt, font=\small},
  % 합산점
  sum/.style   = {draw=ink, line width=0.55pt, fill=white, circle,
                  inner sep=0pt, minimum size=4.6mm, font=\scriptsize},
  asum/.style  = {draw=accent, line width=0.55pt, fill=white, circle,
                  inner sep=0pt, minimum size=4.6mm, font=\scriptsize},
  % 분기점
  dot/.style   = {ink, fill=ink, circle, inner sep=0pt, minimum size=1.5mm},
  % 신호 이름 — 선 위에 작게
  sname/.style = {font=\small, inner sep=1.5pt},
  % 그림 안의 짧은 설명. 그림 안의 글은 최소로 둔다
  note/.style  = {font=\scriptsize, soft, inner sep=1.5pt},
  anote/.style = {font=\scriptsize, accent, inner sep=1.5pt},
  % 축
  axis/.style  = {ink, line width=0.5pt},
  tick/.style  = {font=\scriptsize, soft},
}

% 캡션 한 줄 — 그림 아래 가는 선과 함께
\newcommand{\gnccaption}[2]{%
  \node[anchor=north west, text width=#1, align=left, font=\scriptsize, ink]
        at (current bounding box.south west) {\vspace{2pt}#2};}
  (figures/src/ 안에서)

\usepackage{tikz}
\usepackage{amsmath}
\usepackage{xcolor}
\usetikzlibrary{arrows.meta, positioning, calc, fit, backgrounds, decorations.pathmorphing}

% ---- 색 --------------------------------------------------------------------
% 색은 강조 하나에만 쓴다. 나머지는 검정.
\definecolor{ink}{HTML}{1A1A1A}
\definecolor{accent}{HTML}{7C3AED}
\definecolor{warn}{HTML}{B91C1C}
\definecolor{soft}{HTML}{666666}

% ---- 선과 화살촉 -----------------------------------------------------------
\tikzset{
  % 신호선. 화살촉은 작고 뾰족한 것 하나뿐이다.
  sig/.style   = {ink, line width=0.5pt, -{Stealth[length=4pt, width=3pt]}},
  wire/.style  = {ink, line width=0.5pt},
  % 강조 경로 — 파선으로 구분한다. 흑백 인쇄에서도 살아야 하기 때문이다.
  asig/.style  = {accent, line width=0.5pt, densely dashed,
                  -{Stealth[length=4pt, width=3pt]}},
  awire/.style = {accent, line width=0.5pt, densely dashed},
  % 블록. 흰 바탕, 직각, 검은 테두리.
  blk/.style   = {draw=ink, line width=0.55pt, fill=white,
                  minimum width=17mm, minimum height=8mm, inner sep=2pt, font=\small},
  % 합산점
  sum/.style   = {draw=ink, line width=0.55pt, fill=white, circle,
                  inner sep=0pt, minimum size=4.6mm, font=\scriptsize},
  asum/.style  = {draw=accent, line width=0.55pt, fill=white, circle,
                  inner sep=0pt, minimum size=4.6mm, font=\scriptsize},
  % 분기점
  dot/.style   = {ink, fill=ink, circle, inner sep=0pt, minimum size=1.5mm},
  % 신호 이름 — 선 위에 작게
  sname/.style = {font=\small, inner sep=1.5pt},
  % 그림 안의 짧은 설명. 그림 안의 글은 최소로 둔다
  note/.style  = {font=\scriptsize, soft, inner sep=1.5pt},
  anote/.style = {font=\scriptsize, accent, inner sep=1.5pt},
  % 축
  axis/.style  = {ink, line width=0.5pt},
  tick/.style  = {font=\scriptsize, soft},
}

% 캡션 한 줄 — 그림 아래 가는 선과 함께
\newcommand{\gnccaption}[2]{%
  \node[anchor=north west, text width=#1, align=left, font=\scriptsize, ink]
        at (current bounding box.south west) {\vspace{2pt}#2};}


\begin{document}
\begin{tikzpicture}[x=1mm, y=1mm]

\node[font=\small\bfseries, anchor=north west] at (-34,56)
     {Two headings $20^\circ$ apart, and the error that says $-340^\circ$};
\node[anchor=north west, align=left, text width=150mm, font=\scriptsize, soft] at (-34,50) {%
  Subtracting one angle from another is arithmetic, not geometry. The answer is
  right as a number and wrong as an error, and the vessel turns the long way round.};

% ================= 왼쪽 — 나침반 =================
% 두 방위 사이의 쐐기. 화면 -80 에서 -100 까지 (시계방향 20 deg)
\fill[accent, opacity=0.13] (0,0) -- (-80.00:26.00) arc (-80.00:-100.00:26.00) -- cycle;
\draw[soft, line width=0.4pt] (0,0) circle (26.00);
\draw[axis, -{Stealth[length=4.5pt,width=3.4pt]}] (0,0) -- (0,32.00);
\node[font=\scriptsize, anchor=south] at (0,32.40) {$x^n$~(North)};
\draw[soft, line width=0.4pt] (0.00,26.00) -- (0.00,27.80);
\node[tick] at (30.20,0.00) {E};
\draw[soft, line width=0.4pt] (26.00,0.00) -- (27.80,0.00);
\draw[soft, line width=0.4pt] (0.00,-26.00) -- (0.00,-27.80);
\node[tick] at (-30.20,-0.00) {W};
\draw[soft, line width=0.4pt] (-26.00,-0.00) -- (-27.80,-0.00);

% 긴 호: 화면 -80 -> +260, 반시계 (좌현 340 deg)
\draw[soft, line width=1.5pt, -{Stealth[length=5pt,width=3.8pt]}]
      (-80.00:16.12) arc (-80.00:260.00:16.12);
\node[note, anchor=east, align=right] at (-3.00,16.12)
     {$\psi_d-\psi = -340^\circ$,\\ the long way};

% 짧은 호: 화면 -80 -> -100, 시계 (우현 20 deg)
\draw[accent, line width=1.6pt, -{Stealth[length=5pt,width=3.8pt]}]
      (-80.00:23.40) arc (-80.00:-100.00:23.40);

% 두 방위
\draw[ink, line width=1.0pt, -{Stealth[length=5pt,width=3.6pt]}]
      (0,0) -- (-80.00:26.00);
\node[ink, font=\scriptsize, anchor=west] at (-80.00:29.00)
     {$\psi = +170^\circ$};
\draw[ink, line width=1.0pt, densely dashed, -{Stealth[length=5pt,width=3.6pt]}]
      (0,0) -- (260.00:26.00);
\node[ink, font=\scriptsize, anchor=east] at (260.00:29.00)
     {$\psi_d = -170^\circ$};
\node[anote, anchor=north, align=center, text width=62mm] at (0,-35.50)
     {$\mathrm{ssa}(-340^\circ) = +20^\circ$ — the short way,\\ and the shaded wedge is how wide $20^\circ$ really is};

% ================= 오른쪽 — ssa 는 톱니다 =================
\draw[axis, -{Stealth[length=4pt,width=3pt]}] (55.30,0) -- (122.70,0);
\draw[axis, -{Stealth[length=4pt,width=3pt]}] (88.00,-19.30) -- (88.00,20.30);
\node[font=\scriptsize, anchor=west]  at (123.70,0) {$a$};
\node[font=\scriptsize, anchor=south] at (88.00,20.80) {$\mathrm{ssa}(a)$};
\draw[soft, line width=0.4pt] (68.20,-1.2) -- (68.20,1.2);
\node[tick, anchor=north] at (68.20,-1.6) {$-360$};
\draw[soft, line width=0.4pt] (78.10,-1.2) -- (78.10,1.2);
\node[tick, anchor=north] at (78.10,-1.6) {$-180$};
\draw[soft, line width=0.4pt] (97.90,-1.2) -- (97.90,1.2);
\node[tick, anchor=north] at (97.90,-1.6) {$180$};
\draw[soft, line width=0.4pt] (107.80,-1.2) -- (107.80,1.2);
\node[tick, anchor=north] at (107.80,-1.6) {$360$};
\draw[soft, line width=0.4pt] (86.80,-15.30) -- (89.20,-15.30);
\node[tick, anchor=east] at (86.20,-15.30) {$-180$};
\draw[soft, line width=0.4pt] (86.80,15.30) -- (89.20,15.30);
\node[tick, anchor=east] at (86.20,15.30) {$180$};
\draw[ink, line width=0.9pt] (58.30,-15.30) -- (78.10,15.30);
\draw[ink, line width=0.9pt] (78.10,-15.30) -- (97.90,15.30);
\draw[ink, line width=0.9pt] (97.90,-15.30) -- (117.70,15.30);

% 예제:  a = -340  ->  20
\draw[accent, line width=0.5pt, densely dashed] (69.30,0.00) -- (69.30,1.70) -- (88.00,1.70);
\fill[accent] (69.30,1.70) circle (0.9);
\node[anote, anchor=south west] at (70.50,2.30) {$-340 \mapsto +20$};
\node[note, anchor=north, text width=52mm, align=center] at (88.00,-21.30)
     {Every value is mapped into $(-180^\circ,\,180^\circ]$, and on the
      middle segment $\mathrm{ssa}(a)=a$ leaves it alone. The jumps sit at
      $\pm180^\circ$, which is why the \emph{error} is wrapped, never the heading.};

\end{tikzpicture}
\end{document}
